\documentclass[14pt,a4paper,oneside]{extbook}
\usepackage[utf8]{vietnam}
\usepackage{amsmath,amssymb,mathrsfs,fancyhdr,enumerate,multirow,makecell}
\usepackage{tkz-euclide,tikz-3dplot,tikz,tkz-tab}
\usetikzlibrary{arrows.meta,calc,intersections,angles,quotes,backgrounds,shapes.geometric,patterns,through,shadows,positioning}
\usepackage{yhmath}
\usepackage[top=0.7cm, bottom=0.7cm, left=1.0cm, right=1.0cm, headheight=16pt, headsep=3mm, footskip=5mm]{geometry}
\usepackage[hidelinks,unicode]{hyperref}
\usepackage{qrcode}
\usepackage[many]{tcolorbox}
\usepackage{fontawesome}

\pagestyle{fancy}
\fancyhf{}
\fancyhead[L]{\footnotesize\bfseries\color{teal!85!black}\faBook~SỔ TAY KIẾN THỨC TOÁN 4}
\fancyhead[R]{\footnotesize\bfseries\color{blue!85!black}Lớp Toán Thầy Lĩnh -- loptoanthaylinh.vn}
\fancyfoot[C]{\footnotesize\bfseries\color{gray!80!black}-- Sổ tay kiến thức: Bài 14. So sánh các số có nhiều chữ số --}
\renewcommand{\headrulewidth}{0.4pt}

\begin{document}

% --- BANNER TIÊU ĐỀ BÀI HỌC ---
\noindent
\begin{tcolorbox}[
  enhanced,
  colback=teal!8!white,
  colframe=teal!80!black,
  arc=3mm, boxrule=1pt,
  left=3mm, right=3mm, top=2mm, bottom=2mm,
  halign=center
]
  {\Large\bfseries\color{teal!90!black}\faBookmark\ BÀI 14: SO SÁNH CÁC SỐ CÓ NHIỀU CHỮ SỐ}\\[1.5pt]
  {\footnotesize\itshape\color{black!75} Phương pháp so sánh theo số chữ số, so sánh theo từng hàng và sắp xếp thứ tự trên 1 trang A4}
\end{tcolorbox}

\vspace{2mm}

% --- KHỐI 1 & 2: LÝ THUYẾT & KĨ NĂNG ---
\noindent
\begin{minipage}[t]{0.49\textwidth}
  \begin{tcolorbox}[
    enhanced,
    title={\footnotesize\bfseries\color{white}\faSortAmountDesc\ 1. Quy tắc so sánh 2 số},
    colback=blue!2!white,
    colframe=blue!75!black,
    arc=2.5mm, boxrule=0.8pt, top=1.8mm, bottom=1.8mm, left=2mm, right=2mm
  ]
    \footnotesize
    $\bullet$ \textbf{Trường hợp 1: Khác số chữ số}\\
    Số nào có \textbf{nhiều chữ số hơn} thì số đó \textbf{lớn hơn}.\\
    \textit{VD:} $100\,000$ (6 chữ số) $> 99\,999$ (5 chữ số).\\[1mm]
    $\bullet$ \textbf{Trường hợp 2: Cùng số chữ số}\\
    So sánh từng cặp chữ số ở \textbf{cùng một hàng} kể từ \textbf{trái sang phải}:\\
    --- Ở hàng đầu tiên có sự khác nhau: Số nào có chữ số lớn hơn thì số đó lớn hơn.\\
    --- Nếu tất cả các cặp chữ số bằng nhau thì hai số đó \textbf{bằng nhau}.
  \end{tcolorbox}
\end{minipage}%
\hfill
\begin{minipage}[t]{0.49\textwidth}
  \begin{tcolorbox}[
    enhanced,
    title={\footnotesize\bfseries\color{white}\faTasks\ 2. Kĩ năng sắp xếp thứ tự số},
    colback=teal!2!white,
    colframe=teal!75!black,
    arc=2.5mm, boxrule=0.8pt, top=1.8mm, bottom=1.8mm, left=2mm, right=2mm
  ]
    \footnotesize
    Khi sắp xếp một dãy nhiều số:\\
    $\bullet$ \textbf{Bước 1:} Đếm và phân nhóm các số theo số lượng chữ số (nhóm ít chữ số xếp trước nếu từ bé đến lớn).\\
    $\bullet$ \textbf{Bước 2:} Trong cùng một nhóm có số chữ số bằng nhau, so sánh chữ số ở hàng cao nhất (bên trái cùng).\\
    $\bullet$ \textbf{Bước 3:} Xếp thứ tự theo đúng yêu cầu đề bài: \textbf{Từ bé đến lớn} (tăng dần) hoặc \textbf{Từ lớn đến bé} (giảm dần). Gạch bỏ số đã dùng để tránh trùng lặp.
  \end{tcolorbox}
\end{minipage}

\vspace{2mm}

% --- KHỐI 3: SƠ ĐỒ TƯ DUY TIKZ ---
\noindent
\begin{tcolorbox}[
  enhanced,
  title={\footnotesize\bfseries\color{white}\faSitemap\ SƠ ĐỒ TƯ DUY: THUẬT TOÁN SO SÁNH HAI SỐ TỰ NHIÊN BẤT KÌ},
  colback=yellow!2!white,
  colframe=orange!85!black,
  arc=2.5mm, boxrule=0.8pt, top=1.5mm, bottom=1.5mm, left=1mm, right=1mm,
  halign=center
]
  \begin{tikzpicture}[
    xscale=0.92, yscale=0.82,
    root/.style={rectangle, rounded corners=2mm, draw=teal!80!black, fill=teal!12, thick, inner sep=3.5pt, font=\scriptsize\bfseries, align=center},
    step1/.style={rectangle, rounded corners=2mm, draw=blue!75!black, fill=blue!10, thick, inner sep=3pt, font=\scriptsize\bfseries, align=center},
    step2/.style={rectangle, rounded corners=2mm, draw=orange!80!black, fill=orange!12, thick, inner sep=3pt, font=\scriptsize\bfseries, align=center},
    branch/.style={rectangle, rounded corners=2mm, draw=purple!80!black, fill=purple!10, thick, inner sep=3pt, font=\scriptsize, align=center},
    arr/.style={-{Stealth[scale=0.8]}, thick, draw=black!70}
  ]
    \node[root] (start) at (0,0) {So sánh\\hai số $A$ và $B$};
    \node[step1] (count) at (3.8,0) {Bước 1:\\Đếm số chữ số};

    \node[branch] (diff) at (8.5, 1.2) {\textbf{Khác số chữ số}:\\Số nhiều chữ số hơn $\rightarrow$ \textbf{Lớn hơn}};
    \node[step2] (same) at (8.5, -1.2) {\textbf{Cùng số chữ số}:\\So sánh từ trái sang phải};

    \node[branch] (cmp) at (13.5, -1.2) {Cặp đầu tiên khác nhau:\\Chữ số lớn hơn $\rightarrow$ \textbf{Lớn hơn}};

    \draw[arr] (start) -- (count);
    \draw[arr] (count) |- (diff.west) node[above left, pos=0.8, font=\tiny\bfseries]{Khác nhau};
    \draw[arr] (count) |- (same.west) node[below left, pos=0.8, font=\tiny\bfseries]{Bằng nhau};
    \draw[arr] (same) -- (cmp);
  \end{tikzpicture}
\end{tcolorbox}

\vspace{2mm}

% --- KHỐI 4: VÍ DỤ MINH HỌA & BẪY TRẮC NGHIỆM ---
\noindent
\begin{minipage}[t]{0.53\textwidth}
  \begin{tcolorbox}[
    enhanced,
    title={\footnotesize\bfseries\color{white}\faCheckCircle\ Ví dụ thực hành chuẩn mực},
    colback=green!2!white,
    colframe=green!60!black,
    arc=2.5mm, boxrule=0.8pt, top=1.8mm, bottom=1.8mm, left=2.2mm, right=2.2mm
  ]
    \footnotesize
    $\bullet$ \textbf{Ví dụ 1:} So sánh $89\,999$ và $102\,345$:\\
    Số $89\,999$ có 5 chữ số; $102\,345$ có 6 chữ số.\\
    Vì $5 < 6$ nên $89\,999 < 102\,345$.\\
    $\bullet$ \textbf{Ví dụ 2:} So sánh $745\,280$ và $745\,199$:\\
    Hai số đều có 6 chữ số. So sánh từ trái sang phải:\\
    Hàng trăm nghìn: $7 = 7$; hàng chục nghìn: $4 = 4$; hàng nghìn: $5 = 5$.\\
    Đến hàng trăm: $2 > 1 \Rightarrow \mathbf{745\,280 > 745\,199}$.
  \end{tcolorbox}
\end{minipage}%
\hfill
\begin{minipage}[t]{0.45\textwidth}
  \begin{tcolorbox}[
    enhanced,
    title={\footnotesize\bfseries\color{white}\faExclamationTriangle\ Mẹo tránh bẫy trắc nghiệm},
    colback=red!2!white,
    colframe=red!70!black,
    arc=2.5mm, boxrule=0.8pt, top=1.8mm, bottom=1.8mm, left=2.2mm, right=2.2mm
  ]
    \footnotesize
    $\bullet$ \textbf{Bẫy nhìn trực quan chữ số 9:}\\
    Nhiều học sinh thấy $99\,888$ toàn chữ số 9 lớn nên nghĩ lớn hơn $100\,002$. Hãy luôn \textbf{đếm số chữ số trước}!\\
    $\bullet$ \textbf{Bẫy chiều sắp xếp:}\\
    Đọc kĩ đề yêu cầu "từ bé đến lớn" hay "từ lớn đến bé". Dấu $<$ đi với từ bé đến lớn; dấu $>$ đi với từ lớn đến bé.
  \end{tcolorbox}
\end{minipage}

\end{document}
